    \subsection{FRM EXAM 2009---QUESTION 5-8: ES (99\%),}

    \textbf{Enunciado}

    \emph{Greg Lawrence is a risk analyst at ES Bank. After
    estimating the 99\%, one day VAR of the bank’s portfolio using historical simulation with
    1,200 past days, he is concerned that the VAR measure is not providing enough information
    about tail losses. He decides to reexamine the simulation results. Sorting the simulated daily
    P\&L from worst to best gives the following results:}

    \[
    \begin{array}{c c}
    \toprule
    \text{Rank} & \text{P\&L}\\
    \midrule
    1 & -2{,}833\\
    2 & -2{,}333\\
    3 & -2{,}228\\
    4 & -2{,}084\\
    5 & -1{,}960\\
    6 & -1{,}751\\
    7 & -1{,}679\\
    8 & -1{,}558\\
    9 & -1{,}542\\
    10 & -1{,}484\\
    11 & -1{,}450\\
    12 & -1{,}428\\
    13 & -1{,}368\\
    14 & -1{,}347\\
    15 & -1{,}319\\
    \bottomrule
    \end{array}
    \]


    \emph{What is the 99\%, one-day expected shortfall (ES) of the portfolio?}


    \begin{enumerate}[label=\alph*.]
    \item USD 433
    \item USD 1,428
    \item \hl{USD 1,861}
    \item USD 2,259
    \end{enumerate}

    \subsubsection*{Solución}

    Con \(N=1{,}200\) observaciones, el 1\% peor corresponde a $0.01\times 1{,}200=12 \text{ observaciones.}$
    El ES(99\%) se calcula como el promedio de los peores 12 P\&L.

    Suma:
    \[
    \sum_{i=1}^{12} \text{P\&L}_i=-22330.
    \]
    Promedio:
    \[
    \ES_{0.99}=\frac{-22330}{12}\approx -1860.83.
    \]
    Reportando como magnitud de pérdida:
    \[
    |\ES_{0.99}|\approx 1861.
    \]

    
\paragraph{Comentario.}
Si se promediaran solo las 11 observaciones estrictamente peores que el VaR,
se obtendría
\[
\frac{-20902}{11}\approx -1900.18,
\]
pero ese valor no corresponde a la convención de simulación histórica empleada
en este tipo de ejercicios FRM.










